\section{Conclusions}\label{sec:p2-conclusions}
% P:p2-conclusions-01
A frozen description of water, \COtwo and inorganic ions can accompany amine-specific fitting across MDEA, AMP, PZ and DEA, but calibration and transfer remain property- and source-dependent.
Residual loading-shape errors persist in the fitted pressures, and PZ species residuals remain large relative to their assumed measurement scales.
Shared-parameter substitutions followed by refitting demonstrate compensation: an apparently improved unfitted AMP score before refitting does not survive restoration of its calibration fit.

% P:p2-conclusions-02
MDEA's outside-fit pressure AARDs are 28.19\,\% for Xiao and 18.27\,\% for Ermatchkov, with the latter mixing measured-range and extrapolated dissociation conditions.
AMP's held-out Kundu error is 42.85\,\%, while PZ gives 21.94\,\% against held-out Ermatchkov pressures and 48.52\,\% at 393.15 K.
DEA gives 79.54 and 44.17\,\% on the held-out Seo and Hong and Lal data, despite 0.98\,\% fitted integral-heat AARD, and predicts the wrong sign of water--DEA excess enthalpy.
These unmatched errors are not an amine ranking.

% P:p2-conclusions-03
Local conditional parameter bands are much smaller than the AMP and DEA unfitted discrepancies, overlap the MDEA errors, and are limited by large PZ sensitivities.
At least four false liquid splits remain for the reported PZ fit; a separate penalty fit and direct split evaluations demonstrate that the finite phase search can miss instability.
No additional phase was detected in the specified MDEA, AMP and DEA searches, which is numerical evidence at finite states rather than proof of stability.
The extent of transfer is therefore defined by the reported calibration errors, unfitted property failures, reaction-law extrapolations and conditional uncertainty, rather than by a universal accuracy or process-use claim.
